\section{Conclusion}
\label{sec:conclusion}

This paper presented Cstamp-PWAL, a dependency-closed parallel write-ahead
logging protocol for timestamp-based transaction processing engines.

The core contribution is replacing global durable-prefix acknowledgment with
dependency-closed durable acknowledgment, while reusing the existing commit
timestamp as the logical recovery order.  This design achieves three goals:

\begin{enumerate}

\item \textbf{Preserves recovery safety}: By maintaining a sound over-
approximation of transitive dependencies and ensuring all dependencies are
durable before acknowledging a transaction, Cstamp-PWAL guarantees that crashes
can be recovered correctly without data loss or inconsistency.

\item \textbf{Avoids unnecessary global-prefix coupling}: Transactions are no
longer forced to wait for unrelated WAL shards to flush, eliminating the
conservative durable-acknowledgment backlog that global-prefix systems impose.

\item \textbf{Reduces durable-ack latency on sparse-dependency workloads}: On
YCSB-B with 32 threads, Cstamp-PWAL reduces p99 durable-acknowledgment latency
from 131 milliseconds to 512 microseconds, and reduces pending-commit backlog
from 65K to 55 transactions.  These improvements are critical for latency-
sensitive applications.

\end{enumerate}

The tradeoff is a modest reduction in peak acknowledgment throughput on
unbounded workloads.  However, this is acceptable for systems where latency and
commit backlog are important, and the throughput reduction occurs only in
regimes where global-prefix waiting is not the bottleneck.

\subsection{Future Work}

\subsubsection{Full Crash Recovery Implementation}

The current validation relies on deterministic correctness tests.  A complete
crash recovery implementation would strengthen the paper, especially for top-
tier conference submission.  The design is amenable to full crash recovery
without architectural changes.

\subsubsection{Broader Workload Evaluation}

We evaluated primarily on YCSB-B.  Future work should include TPC-C and other
workloads with varying dependency densities to understand Cstamp-PWAL's
performance across a wider range of transaction patterns.

\subsubsection{Dependency Frontier Optimization}

The current frontier computation incurs approximately 145 microseconds per
transaction in metadata operations.  Opportunities exist to reduce this cost
through better data structures, lock-free algorithms, or compiler optimizations.

\subsubsection{Integration with Other Engines}

While Cstamp-PWAL is demonstrated on ERMIA, the core ideas (dependency-closed
acknowledgment, timestamp reuse as logical LSN) should apply to other timestamp-
based engines like TiKV, CockroachDB, or PostgreSQL with minor modifications.

\subsection{Final Remarks}

Cstamp-PWAL demonstrates that dependency-based durability, when carefully
integrated with timestamp-based transaction engines and evaluated for practical
latency impact, can provide significant improvements over global-prefix designs.
The key insight is that on modern high-concurrency systems, what matters is not
just peak throughput, but also the ability to acknowledge transactions quickly
without accumulating large backlogs.

We hope this work encourages the community to move beyond throughput-only
evaluation and consider latency and backlog as first-class metrics in durability
protocol design.

